\documentclass[12pt,a4paper]{amsart}

\usepackage{amsmath,amssymb,amsthm}
\usepackage{mathtools}
\usepackage[utf8]{inputenc}
\usepackage[T1]{fontenc}
\usepackage{hyperref}
\usepackage{enumitem,booktabs,array}
\usepackage{geometry}
\geometry{margin=2.5cm}

% -------------------------------------------------------
% Theorem environments
% -------------------------------------------------------
\newtheorem{theorem}{Theorem}[section]
\newtheorem{lemma}[theorem]{Lemma}
\newtheorem{corollary}[theorem]{Corollary}
\newtheorem{proposition}[theorem]{Proposition}
\theoremstyle{definition}
\newtheorem{definition}[theorem]{Definition}
\newtheorem{example}[theorem]{Example}
\theoremstyle{remark}
\newtheorem{remark}[theorem]{Remark}
\newtheorem{conjecture}[theorem]{Conjecture}

% -------------------------------------------------------
% Custom commands
% -------------------------------------------------------
\newcommand{\N}{\mathbb{N}}
\newcommand{\Z}{\mathbb{Z}}
\newcommand{\Q}{\mathbb{Q}}
\newcommand{\R}{\mathbb{R}}
\newcommand{\C}{\mathbb{C}}
\newcommand{\Ric}{\mathrm{Ric}}
\newcommand{\Rm}{\mathrm{Rm}}
\newcommand{\vol}{\mathrm{vol}}
\newcommand{\diam}{\mathrm{diam}}
\newcommand{\inj}{\mathrm{inj}}
\newcommand{\Sec}{\mathrm{Sec}}
\DeclareMathOperator{\tr}{tr}
\DeclareMathOperator{\divop}{div}

% -------------------------------------------------------
\begin{document}
% -------------------------------------------------------

\title{Ricci Flow and the Poincar\'{e} Conjecture:\\
       Hamilton's Programme, Perelman's Proof,\\
       and the Geometrization Theorem}

\author{Michael Fuhrmann}
\address{Specialist Mathematics Project, Build 147}
\email{specialist-maths@localhost}
\date{\today}

\begin{abstract}
The Ricci flow, introduced by Richard Hamilton in 1982, is a geometric
evolution equation
\[
  \frac{\partial g}{\partial t} = -2\,\Ric(g)
\]
that deforms a Riemannian metric in the direction of its Ricci curvature.
Hamilton used it to prove that every compact three-manifold with positive
Ricci curvature is diffeomorphic to a spherical space form, and launched a
programme to prove Thurston's Geometrization Conjecture.  In 2002--2003
Grigori Perelman completed this programme by introducing the $\mathcal{F}$-
and $\mathcal{W}$-entropy functionals, the reduced volume, and the technique
of Ricci flow with surgery.  As a corollary, he obtained the Poincar\'{e}
Conjecture: every simply connected, closed, orientable three-manifold is
homeomorphic to $S^3$.  Both results are now fully proved theorems.  We
survey the main definitions, key estimates, the surgery procedure, and the
logical structure of Perelman's argument, and contrast the three-dimensional
situation with the four-dimensional case, where the smooth Poincar\'{e}
Conjecture remains open.
\end{abstract}

\maketitle
\tableofcontents

% ================================================================
\section{Introduction}
% ================================================================

The Poincar\'{e} Conjecture, formulated by Henri Poincar\'{e} in 1904, asks
whether every simply connected closed $3$-manifold is homeomorphic to the
three-sphere $S^3$.  For nearly a century it resisted all attempts at proof
and became the most famous open problem in topology.

Richard Hamilton~\cite{Hamilton1982} proposed in 1982 a strategy based on
\emph{geometric analysis}: deform an arbitrary Riemannian metric by the
\emph{Ricci flow} until it converges to a round metric.  Over the next two
decades he developed deep estimates, classified singularities, and proved
several foundational theorems.  However, the problem of singularities in
dimensions three and above seemed to block a complete proof.

In October 2002 Grigori Perelman posted the first of three preprints on the
arXiv~\cite{Perelman2002,Perelman2003a,Perelman2003b} that completed
Hamilton's programme.  Perelman introduced entirely new monotone quantities
($\mathcal{F}$-entropy, $\mathcal{W}$-entropy, reduced volume), proved the
$\kappa$-non-collapsing theorem, and constructed a canonical Ricci flow with
surgery that continues through all finite-time singularities.  From this he
deduced both the \emph{Geometrization Conjecture} (Thurston 1982) and, as a
special case, the Poincar\'{e} Conjecture.

Perelman was awarded the Fields Medal in 2006 for this work; he declined to
accept it.  Complete, peer-reviewed accounts were subsequently provided by
Morgan--Tian~\cite{MorganTian2007} and Cao--Zhu~\cite{CaoZhu2006}, among
others.

\subsection*{Organisation of this paper}
Section~\ref{sec:riem} reviews the Riemannian geometry background.
Section~\ref{sec:rf} defines the Ricci flow and the normalised Ricci flow.
Section~\ref{sec:hamilton} presents Hamilton's key results.
Section~\ref{sec:surgery} introduces Ricci flow with surgery.
Section~\ref{sec:perelman} describes Perelman's entropy functionals and
monotonicity results.
Section~\ref{sec:poincare} states and proves (by reference) the Poincar\'{e}
theorem.
Section~\ref{sec:geom} treats the Geometrization theorem.
Section~\ref{sec:4d} discusses the four-dimensional situation.
Section~\ref{sec:applications} mentions further applications.


% ================================================================
\section{Riemannian Geometry Background}
\label{sec:riem}
% ================================================================

Let $(M,g)$ be a smooth Riemannian manifold of dimension $n$.

\begin{definition}[Riemann curvature tensor]
The \emph{Riemann curvature tensor} $\Rm$ is the $(1,3)$-tensor defined by
\[
  \Rm(X,Y)Z = \nabla_X\nabla_Y Z - \nabla_Y\nabla_X Z - \nabla_{[X,Y]}Z,
\]
where $\nabla$ is the Levi-Civita connection.  In local coordinates,
$R^l{}_{ijk}$ denotes its components.
\end{definition}

\begin{definition}[Ricci tensor]
The \emph{Ricci tensor} is the contraction
\[
  R_{ij} = R^k{}_{ikj} = g^{kl}R_{kilj}.
\]
It is a symmetric $(0,2)$-tensor measuring, roughly, how much the volume of a
small geodesic ball differs from the Euclidean volume.
\end{definition}

\begin{definition}[Scalar curvature]
The \emph{scalar curvature} is the full trace
\[
  R = g^{ij}R_{ij} = \tr_g \Ric.
\]
\end{definition}

\begin{definition}[Sectional curvature]
For a two-plane $\sigma = \mathrm{span}\{X,Y\}$ in $T_pM$,
\[
  K(\sigma) = \frac{g(\Rm(X,Y)Y,X)}{g(X,X)g(Y,Y)-g(X,Y)^2}.
\]
\end{definition}

\begin{definition}[Einstein manifold]
$(M,g)$ is called \emph{Einstein} if $\Ric = \lambda g$ for some constant
$\lambda \in \R$.  On a compact Einstein three-manifold one has $\lambda > 0$,
$= 0$, or $< 0$ corresponding to positive, zero, or negative scalar curvature.
\end{definition}

\begin{remark}
In dimension three, the full Riemann curvature tensor is completely determined
by the Ricci tensor via
\[
  R_{ijkl} = g_{ik}R_{jl} - g_{il}R_{jk} - g_{jk}R_{il} + g_{jl}R_{ik}
             - \tfrac{1}{2}R(g_{ik}g_{jl} - g_{il}g_{jk}).
\]
This algebraic fact plays a key role in Hamilton's analysis.
\end{remark}

\begin{definition}[Gradient, Laplacian, Divergence]
For a function $f$ and a $(0,2)$-tensor $T$:
\begin{align*}
  (\nabla f)^i &= g^{ij}\partial_j f, \\
  \Delta f &= g^{ij}\nabla_i\nabla_j f, \\
  (\divop T)_j &= g^{ik}\nabla_i T_{kj}.
\end{align*}
\end{definition}


% ================================================================
\section{The Ricci Flow}
\label{sec:rf}
% ================================================================

\begin{definition}[Ricci flow]
Let $M$ be a compact smooth manifold and $g_0$ a Riemannian metric on $M$.
The \emph{Ricci flow} is the one-parameter family of metrics $g(t)$ satisfying
\begin{equation}
  \label{eq:rf}
  \frac{\partial}{\partial t} g(t) = -2\,\Ric(g(t)), \qquad g(0) = g_0.
\end{equation}
\end{definition}

The factor $2$ is conventional; it ensures that the flow coincides with the
heat equation on flat space.

\begin{remark}[Heuristic]
The Ricci flow is a kind of non-linear heat equation for the metric.  Where
curvature is positive (the manifold is ``too curved''), the metric contracts;
where curvature is negative, it expands.  In this way, the flow tends to
homogenise the curvature.
\end{remark}

\begin{theorem}[Short-time existence and uniqueness; DeTurck 1983]
\label{thm:ste}
For every smooth compact Riemannian manifold $(M,g_0)$, there exists
$T > 0$ and a unique smooth solution $g(t)$ of~\eqref{eq:rf} on $[0,T)$.
\end{theorem}

\begin{proof}[Sketch of proof]
The Ricci flow equation is not strictly parabolic due to diffeomorphism
invariance.  DeTurck's trick consists of coupling~\eqref{eq:rf} with a
time-dependent family of diffeomorphisms $\varphi_t$ chosen so that the
pulled-back flow $\tilde g(t) = \varphi_t^* g(t)$ satisfies a strictly
parabolic equation.  Standard parabolic PDE theory then gives short-time
existence and uniqueness of $\tilde g(t)$, from which one recovers $g(t)$.
\end{proof}

\subsection{Evolution equations}

Under the Ricci flow, the geometric quantities evolve as follows:
\begin{align}
  \frac{\partial}{\partial t} R_{ij} &= \Delta R_{ij} + 2 R_{kijl}R^{kl}
     - 2 R_{ik}R^k{}_j + \nabla_i\nabla_j R, \label{eq:ric-ev} \\
  \frac{\partial}{\partial t} R &= \Delta R + 2|\Ric|^2. \label{eq:scal-ev}
\end{align}
Equation~\eqref{eq:scal-ev} implies, via the maximum principle, that the
scalar curvature satisfies $R(x,t) \geq R_{\min}(0)$ for all $t \geq 0$, so
the infimum of $R$ is non-decreasing.

\subsection{Normalised Ricci flow}

To study convergence on manifolds with non-zero Euler characteristic it is
convenient to rescale the metric to keep the volume fixed.  Let
$r(t) = \frac{\int_M R\,d\mu}{\vol(M,g(t))}$ be the average scalar curvature.
The \emph{normalised Ricci flow} is
\begin{equation}
  \label{eq:nrf}
  \frac{\partial}{\partial t} g = -2\Ric + \frac{2}{n} r\, g.
\end{equation}
In dimension $3$ the normalised flow is equivalent to the unnormalised flow up
to a time reparametrisation and a homothetic rescaling of the metric.

\subsection{Ricci solitons}

\begin{definition}[Ricci soliton]
A Riemannian metric $g$ is a \emph{Ricci soliton} if
\[
  \Ric + \nabla^2 f = \lambda g
\]
for some smooth function $f$ (the \emph{potential}) and constant
$\lambda \in \R$.  The soliton is \emph{shrinking} ($\lambda > 0$),
\emph{steady} ($\lambda = 0$), or \emph{expanding} ($\lambda < 0$).
\end{definition}

A Ricci soliton is a fixed point of the normalised Ricci flow modulo
diffeomorphisms and scaling.  The round sphere $S^n$ with its standard metric
is a shrinking Ricci soliton with $f = \mathrm{const}$ (a Ricci--Einstein
metric).  The Gaussian soliton on $\R^n$ is a steady soliton with
$f(x) = |x|^2/4$.


% ================================================================
\section{Hamilton's Results}
\label{sec:hamilton}
% ================================================================

\begin{theorem}[Hamilton 1982]
\label{thm:hamilton82}
Let $(M^3, g_0)$ be a compact three-manifold with positive Ricci curvature,
$\Ric(g_0) > 0$.  Then the normalised Ricci flow exists for all time, and the
metric $g(t)$ converges in $C^\infty$ to a metric of constant positive
sectional curvature as $t\to\infty$.  In particular, $M$ is diffeomorphic to a
spherical space form $S^3/\Gamma$.
\end{theorem}

\begin{proof}[Strategy]
Hamilton analysed the evolution of the curvature tensor using the maximum
principle.  The key steps are:
\begin{enumerate}[label=(\roman*)]
  \item Show that $\Ric > 0$ is preserved under the flow (pinching lemma).
  \item Establish the \emph{pinching estimate}: if $\Ric \geq \delta R\, g$ at
    $t = 0$, then this lower bound improves over time.
  \item Use the Harnack inequality for the curvature to control spatial
    variations.
  \item Deduce exponential convergence to a constant-curvature metric via
    the normalised flow.
\end{enumerate}
The positive Ricci curvature assumption is essential; it prevents the
formation of the long thin ``necks'' ($S^2 \times \R$) that would otherwise
obstruct convergence.
\end{proof}

\begin{corollary}
Every compact three-manifold with positive Ricci curvature has finite
fundamental group.
\end{corollary}

\subsection{Hamilton's entropy}

Hamilton introduced the \emph{entropy functional}
\[
  \mathcal{E}(g) = \int_M R \log R\, d\mu
\]
and showed it is monotone non-decreasing along the normalised Ricci flow on
surfaces, playing a role analogous to Perelman's $\mathcal{W}$-entropy
(Section~\ref{sec:perelman}).

\subsection{Singularity analysis}

In dimensions $\geq 3$, the Ricci flow may develop finite-time singularities.
Hamilton classified them into:
\begin{itemize}
  \item \textbf{Type~I}: $|\Rm|(x,t) \leq \frac{C}{T-t}$ as $t \nearrow T$.
    These include round sphere shrinkage.
  \item \textbf{Type~II}: $|\Rm|$ grows faster than $(T-t)^{-1}$.
    These include the cigar soliton.
  \item \textbf{Type~III}: singularities at $t \to +\infty$ on immortal
    solutions.
\end{itemize}
For three-manifolds with positive Ricci curvature, only Type~I singularities
(corresponding to round collapse) occur, and these are resolved by the
convergence theorem~\ref{thm:hamilton82}.  In the general case without
positivity of Ricci curvature, neck-pinch singularities of Type~I can form;
these require surgery.


% ================================================================
\section{Ricci Flow with Surgery}
\label{sec:surgery}
% ================================================================

\subsection{Neck-pinch singularities}

The most important finite-time singularity in dimension three is the
\emph{neck-pinch}: a region diffeomorphic to $S^2 \times (-1,1)$ develops
unbounded curvature as the $S^2$-factor collapses to a point.

\begin{definition}[$\varepsilon$-neck]
An \emph{$\varepsilon$-neck} in $(M,g)$ centred at $x$ is a region
$U \subset M$ with a diffeomorphism $\psi: S^2 \times (-\varepsilon^{-1},
\varepsilon^{-1}) \to U$ such that
\[
  \bigl\|\, r^{-2}\,\psi^*g - g_{S^2} \oplus ds^2 \,\bigr\|_{C^{[\varepsilon^{-1}]}} < \varepsilon,
\]
where $r$ is the radius of the $S^2$ factor.
\end{definition}

\subsection{Canonical neighbourhood theorem}

Perelman's \emph{canonical neighbourhood theorem} states that any
$(x,t)$ with sufficiently large curvature $|{\Rm}|(x,t) \geq r^{-2}$
has a neighbourhood (after rescaling by $|{\Rm}|(x,t)$) that is
$\varepsilon$-close to a standard model:
\begin{itemize}
  \item a strong $\varepsilon$-neck,
  \item an $\varepsilon$-cap (topology of a half-space or projective plane
    cap), or
  \item a compact $\varepsilon$-round manifold.
\end{itemize}

\subsection{Surgery procedure}

At a singularity time $T$, one performs surgery on all necks:
\begin{enumerate}
  \item Cut $M$ along the neck $S^2$.
  \item Cap off each resulting boundary $S^2$ with a standard round
    hemisphere (\emph{standard cap}).
  \item Restart the flow with the modified metric $g^+(T)$.
\end{enumerate}

\begin{theorem}[Perelman 2003; existence of surgical solution]
\label{thm:surgery}
There exist constants $r, \delta, \kappa > 0$ depending only on the initial
metric such that the Ricci flow with $(r,\delta)$-cutoff exists for all time
$t \in [0,\infty)$, with only finitely many surgeries on any finite time
interval.
\end{theorem}

A key point is that surgeries are performed at a carefully chosen scale $h$
comparable to the curvature scale $r$, so that the geometric quality of the
metric after surgery is controlled.


% ================================================================
\section{Perelman's Proof: Entropy, Non-Collapsing, and Reduced Volume}
\label{sec:perelman}
% ================================================================

\subsection{The $\mathcal{F}$-entropy}

\begin{definition}[Perelman's $\mathcal{F}$-functional]
Let $f: M \to \R$ be a smooth function with $\int_M e^{-f}\,d\mu = 1$.
Define
\begin{equation}
  \mathcal{F}(g,f) = \int_M \bigl(R + |\nabla f|^2\bigr) e^{-f}\,d\mu.
\end{equation}
\end{definition}

\begin{theorem}[Monotonicity of $\mathcal{F}$; Perelman 2002]
\label{thm:F-mono}
Under the \emph{modified Ricci flow}
\[
  \partial_t g_{ij} = -2 R_{ij}, \qquad
  \partial_t f = -\Delta f - R + |\nabla f|^2,
\]
the $\mathcal{F}$-functional is non-decreasing:
\begin{equation}
  \frac{d}{dt}\mathcal{F}(g(t),f(t))
  = 2\int_M \bigl|R_{ij} + \nabla_i\nabla_j f\bigr|^2 e^{-f}\,d\mu \geq 0.
\end{equation}
Equality holds if and only if $(g,f)$ is a gradient steady Ricci soliton.
\end{theorem}

\begin{proof}[Sketch]
Direct computation using the evolution equations for $g$, $d\mu$, $R$, and
$|\nabla f|^2$.  Integration by parts converts boundary terms to zero on a
compact manifold.
\end{proof}

\subsection{The $\mathcal{W}$-entropy}

\begin{definition}[Perelman's $\mathcal{W}$-functional]
For $\tau > 0$ and $(g, f)$ with $\int_M (4\pi\tau)^{-n/2}e^{-f}\,d\mu = 1$:
\begin{equation}
  \mathcal{W}(g,f,\tau)
  = \int_M \bigl[\tau(R + |\nabla f|^2) + f - n\bigr]
    (4\pi\tau)^{-n/2}e^{-f}\,d\mu.
\end{equation}
\end{definition}

\begin{definition}[Perelman entropy $\mu$-functional]
\[
  \mu(g,\tau) = \inf \bigl\{\mathcal{W}(g,f,\tau) :
    \int_M (4\pi\tau)^{-n/2}e^{-f}\,d\mu = 1 \bigr\}.
\]
\end{definition}

\begin{theorem}[Monotonicity of $\mathcal{W}$; Perelman 2002]
\label{thm:W-mono}
Under the Ricci flow coupled with
\[
  \partial_t f = -\Delta f - R + \frac{n}{2\tau}, \qquad
  \frac{d\tau}{dt} = -1,
\]
the $\mathcal{W}$-functional is non-decreasing.  In particular,
$t \mapsto \mu(g(t), T-t)$ is non-decreasing for $t < T$.
\end{theorem}

\subsection{$\kappa$-non-collapsing}

\begin{theorem}[$\kappa$-non-collapsing; Perelman 2002]
\label{thm:kappa}
Let $(M,g(t))$ be a Ricci flow on $[0,T]$.  For every $\rho > 0$ there
exists $\kappa > 0$ (depending on $\rho$, $T$, and $g(0)$) such that the
following holds: if $|\Rm|(x,t) \leq \rho^{-2}$ on the ball $B(x,t,\rho)$,
then
\[
  \vol_t B(x,t,\rho) \geq \kappa \rho^n.
\]
\end{theorem}

\begin{proof}[Strategy]
The $\mathcal{W}$-entropy lower bound provides a lower bound on
$\mu(g(t),\rho^2)$.  A collapsing ball would force $\mathcal{W}$ to be very
negative, contradicting the monotonicity theorem~\ref{thm:W-mono}.
\end{proof}

The $\kappa$-non-collapsing theorem is the central tool: it prevents the
development of the cigar-type singularities that had previously obstructed
Hamilton's programme, and ensures that high-curvature regions have the
canonical neighbourhood structure needed for surgery.

\subsection{Reduced volume and reduced length}

Perelman defined the \emph{$\mathcal{L}$-geodesic} and \emph{reduced length}:
\begin{equation}
  l(q,\tau) = \frac{1}{2\sqrt{\tau}}
    \inf_\gamma \int_0^\tau \sqrt{s}\bigl(R(\gamma(s)) + |\dot\gamma(s)|^2\bigr)ds,
\end{equation}
where the infimum is over all curves $\gamma$ from $(p,0)$ to $(q,\tau)$
traversed backwards in time.  The \emph{reduced volume} is
\begin{equation}
  \tilde V(\tau) = \int_M (4\pi\tau)^{-n/2} e^{-l(q,\tau)}\,d\vol_{g(\tau)}.
\end{equation}

\begin{theorem}[Monotonicity of reduced volume; Perelman 2002]
Under the Ricci flow, $\tilde V(\tau)$ is non-increasing in $\tau$, and
$\tilde V(\tau) \leq 1$.  If $\tilde V(\tau_0) = 1$ for some $\tau_0 > 0$,
then $(M,g)$ is isometric to Euclidean space.
\end{theorem}


% ================================================================
\section{The Poincar\'{e} Conjecture}
\label{sec:poincare}
% ================================================================

\begin{theorem}[Poincar\'{e} Conjecture; Perelman 2003]
\label{thm:poincare}
Every simply connected, closed (compact, without boundary), orientable
three-manifold is homeomorphic to the three-sphere $S^3$.
\end{theorem}

\begin{proof}[Proof via Ricci flow with surgery]
Let $M$ be a simply connected closed three-manifold.  Equip $M$ with an
arbitrary smooth metric $g_0$ and run the Ricci flow with surgery
(Theorem~\ref{thm:surgery}).

\medskip
\textbf{Step 1: The flow becomes extinct in finite time.}
Since $\pi_1(M) = 0$, in particular $\pi_2(M) = 0$ (by Hurewicz) and the
$3$-manifold is prime.  Perelman~\cite{Perelman2003b} (and independently
Colding--Minicozzi) proved that for simply connected closed $3$-manifolds
the Ricci flow with surgery becomes extinct in finite time $T < \infty$:
the manifold is completely decomposed after finitely many surgeries, each
piece becoming round before collapsing.

\medskip
\textbf{Step 2: No non-trivial topology survives.}
Just before extinction, the metric is $\varepsilon$-close to a round metric
on $S^3$ or on a spherical space form $S^3/\Gamma$.  Since $\pi_1(M) = 0$,
no non-trivial $\Gamma$ can arise, so the only possibility is $S^3$.

\medskip
\textbf{Step 3: Conclusion.}
All surgeries replace $M$ by a topologically equivalent manifold (by the
canonical neighbourhood structure, surgeries are performed along $S^2$-necks
and attach standard caps that are topological balls).  Therefore the original
manifold $M$ is homeomorphic to $S^3$.
\end{proof}

\begin{remark}
The complete verification of Perelman's argument was carried out by
Morgan--Tian~\cite{MorganTian2007} and Cao--Zhu~\cite{CaoZhu2006}.  Both
independently confirmed all details of the proof.  There is now complete
consensus in the mathematical community that Theorem~\ref{thm:poincare} is
established.
\end{remark}

\begin{remark}[Historical context]
The Poincar\'{e} Conjecture was one of the seven Millennium Prize Problems
of the Clay Mathematics Institute (CMI), each carrying a prize of
\$1{,}000{,}000.  Perelman was awarded the prize in 2010 but declined to
accept it, just as he had declined the Fields Medal in 2006.
\end{remark}


% ================================================================
\section{The Geometrization Conjecture}
\label{sec:geom}
% ================================================================

\begin{definition}[Thurston's eight geometries]
\label{def:geom}
A three-dimensional geometry $(X,G)$ consists of a simply connected
three-manifold $X$ together with a transitive isometry group $G$ such that
the stabiliser of a point is compact.  There are exactly eight such maximal
geometries:
\[
  S^3,\quad \R^3,\quad H^3,\quad S^2\times\R,\quad H^2\times\R,\quad
  \widetilde{\mathrm{SL}_2(\R)},\quad \mathrm{Nil},\quad \mathrm{Sol}.
\]
\end{definition}

\begin{definition}[Geometric decomposition]
A closed, orientable three-manifold $M$ is \emph{geometric} if it is modelled
on one of the eight geometries.  The \emph{JSJ decomposition} cuts $M$ along
incompressible tori into atoroidal pieces; together with the prime
decomposition along $S^2$, one obtains the \emph{geometrisation}.
\end{definition}

\begin{theorem}[Geometrization Theorem; Thurston 1982, Perelman 2003]
\label{thm:geom}
Every closed, orientable three-manifold can be cut along finitely many
disjoint, incompressible tori and spheres into pieces, each of which
admits a complete, locally homogeneous Riemannian metric modelled on one of
the eight Thurston geometries.
\end{theorem}

\begin{proof}[Proof via Ricci flow with surgery]
Perelman's Ricci flow with surgery provides a canonical decomposition of $M$
into prime pieces.  For the atoroidal (hyperbolic) pieces, Perelman established
the \emph{thick-thin decomposition}: after rescaling, the thick part converges
to a hyperbolic manifold (using Hamilton's compactness theorem and the
$\kappa$-non-collapsing estimates), while the thin part is a graph manifold
(and hence carries one of the Seifert-fibred geometries).  The graph manifold
case was handled separately using results of Shioya--Yamaguchi.  Together
these give the full Geometrization Theorem.
\end{proof}

\begin{corollary}
The Poincar\'{e} Conjecture (Theorem~\ref{thm:poincare}) follows from the
Geometrization Theorem: the only geometry compatible with a simply connected
closed three-manifold is the spherical geometry $S^3$, so $M \cong S^3$.
\end{corollary}

\begin{remark}
The Geometrization Theorem was first conjectured by William Thurston~\cite{Thurston1982}
in 1982, based on his deep work on hyperbolic three-manifolds (for which he
received the Fields Medal in 1982).  Thurston proved the conjecture for
Haken manifolds; the general case required Perelman's Ricci flow methods.
\end{remark}


% ================================================================
\section{The Four-Dimensional Case}
\label{sec:4d}
% ================================================================

The situation in dimension four differs fundamentally from dimension three.

\subsection{Topological Poincar\'{e} conjecture in dimension 4}

\begin{theorem}[Freedman 1982]
\label{thm:freedman}
Every simply connected closed topological four-manifold with the homotopy
type of $S^4$ is homeomorphic to $S^4$.
\end{theorem}

Freedman's proof~\cite{Freedman1982} uses the technique of
\emph{Casson handles} and is specific to dimension four.  It builds on the
Whitney trick, which works topologically in dimension four but not smoothly.
Freedman was awarded the Fields Medal in 1986 for this work.

\subsection{Smooth Poincar\'{e} conjecture in dimension 4 --- open problem}

\begin{conjecture}[Smooth four-dimensional Poincar\'{e} conjecture]
\label{conj:s4d}
Every smooth manifold homeomorphic to $S^4$ is diffeomorphic to $S^4$.
\end{conjecture}

This remains one of the most important open problems in geometric topology.
In contrast to the topological statement (Theorem~\ref{thm:freedman}), the
smooth problem is completely open.

\begin{remark}
In all other dimensions $n \neq 4$, the smooth Poincar\'{e} conjecture is
decided:
\begin{itemize}
  \item $n \leq 2$: classical (and trivial for $n=1$).
  \item $n = 3$: Perelman's theorem (Theorem~\ref{thm:poincare}).
  \item $n \geq 5$: Smale (1961) proved that every simply connected, closed
    $n$-manifold homotopy equivalent to $S^n$ is diffeomorphic to $S^n$
    for $n \geq 5$ (the $h$-cobordism theorem).
\end{itemize}
Dimension four is exceptional because the Whitney trick works topologically
but not smoothly.  Donaldson gauge theory and Seiberg--Witten theory have
produced obstructions showing that exotic smooth structures exist on $\R^4$
(but not on $\R^n$ for $n \neq 4$), and many simply connected $4$-manifolds
admit infinitely many distinct smooth structures.  For $S^4$ itself, no exotic
smooth structure has been found, but none has been ruled out.
\end{remark}

\subsection{Ricci flow in dimension 4}

The Ricci flow in dimension four faces new challenges:
\begin{itemize}
  \item The Weyl curvature tensor does not vanish in dimension four, so the
    Riemann tensor is not determined by $\Ric$ alone.
  \item New types of singularities (e.g., orbifold singularities) appear.
  \item Hamilton's classification of singularities is less complete.
\end{itemize}
Progress includes work of Chen--Zhu and of Brendle on four-manifolds with
pointwise $\frac{1}{4}$-pinched sectional curvature.


% ================================================================
\section{Applications and Generalisations}
\label{sec:applications}
% ================================================================

\subsection{Mean curvature flow}

The \emph{mean curvature flow} (MCF) moves a hypersurface $\Sigma_t \subset \R^{n+1}$
in the direction of its mean curvature vector:
\[
  \frac{d}{dt}\mathbf{x} = H\,\mathbf{n},
\]
where $H$ is the mean curvature and $\mathbf{n}$ the inward unit normal.
By analogy with Ricci flow, the MCF develops singularities (neck-pinches) that
can be resolved by surgery.  Andrews, Brendle, Huisken, and Ilmanen have
developed a MCF with surgery analogous to Perelman's construction.

\subsection{K\"{a}hler--Ricci flow}

On a K\"{a}hler manifold $(M, J, g)$ the Ricci form $\rho = -\frac{i}{2\pi}\partial\bar\partial\log\det(g_{i\bar j})$
represents the first Chern class.  The \emph{K\"{a}hler--Ricci flow} is
\[
  \frac{\partial \omega}{\partial t} = -\Ric(\omega) + \lambda\omega, \qquad \lambda \in \{-1, 0, 1\},
\]
and is used in the study of canonical K\"{a}hler metrics (Calabi--Yau,
K\"{a}hler--Einstein).  The Yau--Tian--Donaldson conjecture (proved by Chen,
Donaldson, and Sun in 2015) states that K-stability is equivalent to the
existence of a K\"{a}hler--Einstein metric, and the K\"{a}hler--Ricci flow
plays a central role.

\subsection{Harmonic map flow and Yang--Mills flow}

The techniques developed for Ricci flow have analogues in:
\begin{itemize}
  \item \textbf{Harmonic map flow}: $\partial_t u = \tau(u)$, where $\tau$ is
    the tension field.  Eells--Sampson (1964) proved long-time existence when
    the target has non-positive curvature.
  \item \textbf{Yang--Mills flow}: gradient flow of the Yang--Mills energy
    functional.  Donaldson used it to prove the Donaldson--Uhlenbeck--Yau
    theorem on stable bundles.
\end{itemize}

\subsection{Thurston's program and low-dimensional topology}

The Geometrization Theorem resolves the \emph{virtual Haken conjecture} (proved by
Agol 2012) and the \emph{virtual fibering conjecture}, both now theorems.
The classification of three-manifolds is now, in principle, complete: every
closed orientable three-manifold can be understood in terms of the eight
model geometries.


% ================================================================
\section*{Acknowledgements}
% ================================================================

This paper is part of the \emph{Specialist Mathematics Project}, a systematic
study of advanced mathematics from algebra and analysis to the Millennium
Problems.  The exposition follows the treatments in Morgan--Tian~\cite{MorganTian2007},
Cao--Zhu~\cite{CaoZhu2006}, and Kleiner--Lott~\cite{KleinerLott2008}.


% ================================================================
\begin{thebibliography}{99}
% ================================================================

\bibitem{Hamilton1982}
R.~S.~Hamilton,
\textit{Three-manifolds with positive Ricci curvature},
J.\ Differential Geom.\ \textbf{17} (1982), no.~2, 255--306.

\bibitem{Perelman2002}
G.~Perelman,
\textit{The entropy formula for the Ricci flow and its geometric applications},
arXiv:math/0211159 (2002).

\bibitem{Perelman2003a}
G.~Perelman,
\textit{Ricci flow with surgery on three-manifolds},
arXiv:math/0303109 (2003).

\bibitem{Perelman2003b}
G.~Perelman,
\textit{Finite extinction time for the solutions to the Ricci flow on certain three-manifolds},
arXiv:math/0307245 (2003).

\bibitem{MorganTian2007}
J.~Morgan and G.~Tian,
\textit{Ricci Flow and the Poincar\'{e} Conjecture},
Clay Mathematics Monographs, vol.~3, AMS, Providence, 2007.

\bibitem{CaoZhu2006}
H.-D.~Cao and X.-P.~Zhu,
\textit{A complete proof of the Poincar\'{e} and geometrization conjectures ---
application of the Hamilton-Perelman theory of the Ricci flow},
Asian J.\ Math.\ \textbf{10} (2006), no.~2, 165--492.

\bibitem{KleinerLott2008}
B.~Kleiner and J.~Lott,
\textit{Notes on Perelman's papers},
Geom.\ Topol.\ \textbf{12} (2008), 2587--2855.

\bibitem{Thurston1982}
W.~P.~Thurston,
\textit{Three-dimensional manifolds, Kleinian groups and hyperbolic geometry},
Bull.\ Amer.\ Math.\ Soc.\ \textbf{6} (1982), 357--381.

\bibitem{Freedman1982}
M.~H.~Freedman,
\textit{The topology of four-dimensional manifolds},
J.\ Differential Geom.\ \textbf{17} (1982), 357--453.

\bibitem{Chow2004}
B.~Chow and D.~Knopf,
\textit{The Ricci Flow: An Introduction},
Mathematical Surveys and Monographs, vol.~110, AMS, Providence, 2004.

\end{thebibliography}

\end{document}
